\section{二次问题的迭代}

在这一节，我们固定考虑一个特定的二次问题，形如
\begin{Equation}[二次问题的目标函数]
    \begin{array}{ll}
        \minimize_{\vb{x}}&f(\vb{x})=\frac{1}{2}(\vb{x}-\vb{x}^{*})^T\vb{Q}(\vb{x}-\vb{x}^{*})\\
        \subto&\vb{x}\in\R
    \end{array}
\end{Equation}

其中$\vb{Q}\in\R^{n\times n}$并满足$\vb{Q}\succeq\vb{0}$，显然
\begin{Equation}[二次问题的梯度]
    \begin{array}{ll}
       \grad f(\vb{x})=\vb{Q}(\vb{x}-\vb{x}^{*})
    \end{array}
\end{Equation}

接下来，我们将在该二次问题讨论两种迭代步长的选取：固定步长，精确线搜索。

\subsection{固定步长的收敛性}

\begin{BoxProperty}[二次问题在固定步长下的步长]
    二次问题，应用固定步长，应取步长为
    \begin{Equation}[二次问题在固定步长下的步长]
        \eta_t=\frac{2}{\lambda_1+\lambda_n}
    \end{Equation}
    其中$\lambda_1$和$\lambda_n$分别是$\vb{Q}$最大和最小的特征值。
\end{BoxProperty}

\begin{BoxProperty}[二次问题在固定步长下的收敛]
    二次问题，应用固定步长，满足以下收敛性质
    \begin{Equation}[二次问题在固定步长下的收敛]
        \norm{\vb{x}_t-\vb{x}}\leq\qty(\frac{\lambda_1-\lambda_n}{\lambda_1-\lambda_n})^t\norm{\vb{x}_0-\vb{x}^{*}}
    \end{Equation}
\end{BoxProperty}

我们可以证明这一点，根据\cref{def:梯度下降}，代入\cref{eq:二次问题的梯度}
\begin{Equation}[二次问题固定步长1]
    \vb{x}_{t+1}=\vb{x}_t-\eta\grad f(\vb{x}_t)=\vb{x}_t-\eta\vb{Q}(\vb{x}_t-\vb{x}^{*})
\end{Equation}

两边同时减去$\vb{x}^{*}$
\begin{Equation}[二次问题固定步长2]
    \vb{x}_{t+1}-\vb{x}^{*}=(\vb{I}-\eta\vb{Q})(\vb{x}_t-\vb{x}^{*})
\end{Equation}

两边取模
\begin{Equation}[二次问题固定步长2]
    \norm{\vb{x}_{t+1}-\vb{x}^{*}}_2\leq\norm{\vb{I}-\eta\vb{Q}}_2\norm{\vb{x}_t-\vb{x}^{*}}_2
\end{Equation}

矩阵的$\norm{\cdot}_2$是谱范数，谱范数等价于最大奇异值，而对称矩阵的奇异值就是特征值的模
\begin{Equation}[二次问题固定步长3]
    \norm{\vb{I}-\eta\vb{Q}}_2=\max\abs{1-\eta\lambda}
\end{Equation}

显然，问题的关键在于$\lambda_1$和$\lambda_n$
\begin{Equation}[二次问题固定步长4]
    \norm{\vb{I}-\eta\vb{Q}}_2=\max\qty{\abs{1-\eta\lambda_1},\abs{1-\eta\lambda_n}}
\end{Equation}

因此$\eta$的取值应能使两者相同，且由于$\lambda_1>\lambda_n$，两者恰好是一正一负
\begin{Equation}[二次问题固定步长5]
    (1-\eta\lambda_1)=-(1-\eta\lambda_n)
\end{Equation}

从而解得
\begin{Equation}[二次问题固定步长6]
    \eta=\frac{2}{\lambda_1+\lambda_n}
\end{Equation}

以及
\begin{Equation}[二次问题固定步长7]
    \abs{1-\eta\lambda_n}=\abs{1-\eta\lambda_1}=\frac{\lambda_1+\lambda_n}{\lambda_1-\lambda_n}
\end{Equation}

只要逐次应用不等式，就可以得到收敛条件了。

\subsection{精确线搜索的收敛性}

然而，通常不可能事先确定特征值。更实际的方案是精确线搜索（Exact Line Search），它的思想非常简单，每一次迭代的步长$\eta_t$取决于什么样的$\eta_t$可以让$f(\vb{x}_t-\eta_t\grad f(\vb{x}_t))$下降最多？

\begin{BoxDefinition}[精确线搜索]
    精确线搜索是一种确定步长的方法
    \begin{Equation}[精确线搜索]
        \eta_t=\arg\min_{\eta\geq 0}f(\vb{x}_t-\eta\grad f(\vb{x}_t))
    \end{Equation}
\end{BoxDefinition}

在二次问题下，精确线搜索的步长和收敛可以被表示为

\begin{BoxProperty}[二次问题在精确线搜索下的步长]
    二次问题，应用精确线搜索，应取步长为
    \begin{Equation}[二次问题在固定步长下的步长]
        \eta_t=\frac{\vb{g}_t^T\vb{g}_t}{\vb{g}_T^T\vb{Q}\vb{g}_t}
    \end{Equation}
    其中$\vb{g}_t=\grad f(\vb{x}_t)=\vb{Q}(\vb{x}_t-\vb{x}^{*})$。
\end{BoxProperty}

\begin{BoxProperty}[二次问题在精确线搜索下的收敛]
    二次问题，应用精确线搜索，满足以下收敛性质
    \begin{Equation}[二次问题在精确线搜索下的收敛]
        f(\vb{x}_t)-f(\vb{x}^{*})\leq\qty(\frac{\lambda_1-\lambda_n}{\lambda_1+\lambda_n})^{2t}(f(\vb{x}_0-f\vb{x}^{*}))
    \end{Equation}
\end{BoxProperty}

有些情况下，收敛性会用$f(\vb{x}_t)$和$f(\vb{x}^{*})$而不是$\vb{x}_t$和$\vb{x}^{*}$的接近程度来表征。





